\chapter{CDM Duties, Documentation and Competence Risk Assessment}
\label{chap:CDM Duties, Documentation and Competence Risk Assessment}

%---------------------------- SECTION ----------------------------
\section{Why this assessment matters in construction}

\paragraph{The Construction (Design and Management) Regulations 2015 represent the most comprehensive and far-reaching piece of construction-specific health and safety legislation in the United Kingdom. CDM 2015 restructures the allocation of duties across the entire construction project team --- from the client who commissions the work through the designers who shape the hazard profile of the built asset, through the principal contractor who manages the construction phase, to the workers who execute the physical work. Understanding CDM 2015 not as a bureaucratic compliance framework but as a duty-led risk management system is fundamental to the effective assessment and control of construction health and safety risk at every level of the project structure.}

\paragraph{The introduction of CDM 2015 replaced the CDM 2007 regime and the previous CDM Coordinator role with the Principal Designer --- a duty holder who must be an organisation or individual with control over the pre-construction phase of a project and who carries specific obligations to plan, manage, monitor, and coordinate health and safety during the pre-construction phase. This change was intended to embed design-phase risk management into the project structure from the outset, recognising the well-established principle that the majority of construction health and safety risk is determined by decisions made at design stage --- and that those decisions are most effectively addressed during design, not after the design is fixed and construction has commenced.}

\paragraph{CDM documentation --- the Pre-Construction Information, the Construction Phase Plan, and the Health and Safety File --- forms the evidential backbone of the construction project's risk management system. These documents are not administrative formalities; they are the mechanisms through which risk information is gathered, communicated, acted upon, and preserved for the benefit of the building's future maintainers and occupants. The absence of a Construction Phase Plan before construction commences is a criminal offence; the failure to prepare a Health and Safety File and pass it to the client at completion denies future occupants, maintainers, and demolition contractors the information they need to work safely in and on the building.}

\paragraph{Competence --- the knowledge, skills, experience, and behaviours necessary to perform one's CDM duties safely and effectively --- is the thread that runs through every CDM obligation. The shift from CDM 2007's ``competence assessment'' requirement (which had become a box-ticking exercise through the proliferation of pre-qualification schemes) to the CDM 2015 approach of proportionate competence assessment has placed greater emphasis on the subjective judgement of the appointing duty holder. Clients must ensure that those they appoint have the skills, knowledge, experience, and organisational capability to perform their CDM roles; and this assessment must be genuine, not merely the verification that a CSCS card has been issued or an SSIP certificate presented.}

\barquote{``CDM is not a form-filling exercise. It is the structural framework through which the health and safety obligations of every person who touches a construction project are defined, allocated, and enforced. A CDM compliance document that exists without the culture that should underpin it is a legal fiction.''}

\example{CDM Scenario}{A domestic client commissions a builder to construct a large home extension. The project involves groundworks, structural work, first-floor construction, and a new roof. Total construction value is \pounds180\,000 and duration will be 14 weeks. The builder, who is the only contractor, proceeds without preparing a Construction Phase Plan and without notifying the HSE under F10, on the grounds that the client is a domestic client and the project is therefore exempt. The builder is incorrect on both counts: the domestic client provision of CDM 2015 does not exempt the project from CDM --- it transfers the client's CDM duties to the builder as the only contractor. Furthermore, the project exceeds the notification thresholds (more than 30 working days with more than 20 workers simultaneously, or more than 500 person-days). An HSE inspector notified by a concerned structural engineer issues a Prohibition Notice stopping work until an F10 is submitted and a Construction Phase Plan produced.}

%---------------------------- SECTION ----------------------------
\section{Legal and professional framework}

\paragraph{CDM 2015 is a set of Regulations made under the Health and Safety at Work etc.\ Act 1974, supplemented by the L153 Approved Code of Practice and Guidance which provides the definitive interpretation of the Regulations' requirements. L153 has status in law: a failure to comply with L153 is not automatically a criminal offence, but it is admissible as evidence of failure to comply with the underlying Regulation, and the defendant must demonstrate that they complied with the Regulation by other equally effective means. The HSE's suite of CDM guidance documents, including the specific guidance for clients, designers, principal designers, principal contractors, and workers, provides further interpretation of the L153 approach.}

\paragraph{CDM 2015 applies to all construction work in Great Britain. ``Construction work'' is defined broadly in Regulation 2 and includes demolition, maintenance, alteration, conversion, fitting-out, and commissioning of a structure, as well as the preparatory work and groundwork necessary for a structure. The breadth of this definition means that the vast majority of building work, regardless of scale, is subject to CDM 2015. The Regulations do not apply in Northern Ireland, which has its own CDM Regulations modelled on the 2007 regime.}

\begin{itemize}
  \bitem{CDM 2015 --- Client duties (Regulations 4 and 5)}{The client must ensure suitable arrangements for managing the project (including allocation of sufficient time and resources); appoint a Principal Designer in writing for projects with more than one contractor; appoint a Principal Contractor in writing; ensure that the pre-construction information is provided to designers and contractors; ensure that the Construction Phase does not start until suitable welfare arrangements and a Construction Phase Plan are in place; and ensure that the Health and Safety File is prepared, kept, reviewed, and passed to a new client on transfer of ownership.}
  \bitem{CDM 2015 --- Principal Designer duties (Regulations 11 and 12)}{The Principal Designer must plan, manage, monitor, and coordinate health and safety in the pre-construction phase; take account of the General Principles of Prevention in preparing, reviewing, and updating the design; identify, eliminate, or control foreseeable risks in the construction, maintenance, and end-of-life phases; ensure cooperation and communication between designers; prepare and update the pre-construction information; and produce and update the Health and Safety File.}
  \bitem{CDM 2015 --- Principal Contractor duties (Regulations 12 to 14)}{The Principal Contractor must plan, manage, monitor, and coordinate health and safety in the construction phase; prepare, develop, and implement the Construction Phase Plan; organise cooperation between contractors; ensure that reasonable steps are taken to prevent unauthorised access to the site; ensure that all contractors are provided with relevant pre-construction information; ensure that workers have been consulted on matters affecting their health and safety; and liaise with the Principal Designer over the health and safety information to be included in the Health and Safety File.}
  \bitem{Pre-Construction Information (PCI)}{The PCI is the collection of information about the project gathered by or for the client and provided to designers and contractors before and during the design process. It includes: survey and geotechnical information; information about existing structures and services; information about the site environment (contamination, adjacent uses, traffic); and information about hazardous materials in existing structures to be demolished or modified (asbestos, lead paint, PCB-containing electrical equipment). The adequacy of the PCI directly affects the quality of the design risk assessment and the Construction Phase Plan.}
  \bitem{Construction Phase Plan (CPP)}{The CPP is prepared by the Principal Contractor (or sole contractor) before the construction phase starts and is developed throughout the construction phase. It describes the management structure for health and safety on the project; the arrangements for managing the significant risks identified in the PCI and by the Principal Contractor's own risk assessment; site rules; welfare arrangements; emergency procedures; and arrangements for monitoring health and safety performance. The CPP is a living document, not a pre-start form.}
  \bitem{Health and Safety File (HSF)}{The HSF is prepared by the Principal Designer and updated throughout the construction phase by the Principal Contractor, containing information needed for future maintenance and use of the structure: as-built drawings, structural engineering design data, O\&M manuals, record of hazardous materials, and information about embedded services and their isolation arrangements. The HSF must be handed to the client at completion and must be made available to anyone who carries out construction work on the building in future.}
  \bitem{F10 Notification}{Projects that will last more than 30 working days and have more than 20 workers simultaneously, or will exceed 500 person-days, must be notified to the HSE using the F10 form (or its online equivalent). The notification must be made as soon as practicable before the construction phase begins. The notification is the client's duty; it is typically delegated to the Principal Contractor in practice. Failure to notify is a criminal offence.}
  \bitem{CSCS Card Scheme}{The Construction Skills Certification Scheme provides cards to construction workers and supervisors as evidence that they hold appropriate health and safety training and qualifications for their occupation. The CSCS scheme does not confer competence but provides a minimum assurance of health and safety awareness. Principal contractors typically require CSCS cards as a condition of site access; the card level (Green Labourer, Blue Skilled Worker, Gold Supervisor, Black Manager) must be appropriate to the role.}
\end{itemize}

%---------------------------- SECTION ----------------------------
\section{Conducting the assessment using the five-step method}

\subsection{Step 1 --- Identify Hazards}
\paragraph{CDM-related hazards arise from failures in the duty holder structure and documentation framework. Key hazards include: absence of a competent Principal Designer, leading to inadequate design risk assessment and deficient pre-construction information that passes uncontrolled hazards to the construction phase; absence of or deficient Construction Phase Plan, resulting in unmanaged construction phase risks that are foreseeable but unprepared for; absence of competence assessment for appointed duty holders, leading to organisations and individuals performing CDM roles they lack the knowledge and skills to discharge; failure to notify the HSE where notification is required; failure to provide adequate pre-construction information to contractors and designers; failure to complete or pass on the Health and Safety File, denying future maintainers and occupants essential safety information; and domestic client provisions being misapplied to exclude projects from CDM requirements that in fact apply.}

\subsection{Step 2 --- Decide Who or What Is Affected}
\paragraph{All persons who work on, occupy, maintain, or eventually demolish the building are affected by the quality of the CDM documentation. Construction phase workers are directly affected by the adequacy of the Construction Phase Plan; a deficient CPP that fails to address a significant risk (say, the presence of asbestos in a building being refurbished) may result in operatives being exposed to a hazard that the plan should have controlled. Future maintainers are affected by the completeness of the Health and Safety File; a roofer conducting maintenance on a building whose HSF does not record that the roof is not designed for maintenance access may not appreciate the risk of falling through fragile rooflight panels until it is too late. Members of the public adjacent to the construction site are affected by the adequacy of the Construction Phase Plan's site security and public protection provisions.}

\subsection{Step 3 --- Evaluate Likelihood and Consequence}
\paragraph{The consequence of CDM failures can be severe across all three documentation domains. A deficient PCI that fails to identify asbestos can result in illegal asbestos exposure --- a risk of mesothelioma that may not manifest for 30 to 40 years. A deficient CPP that fails to address the support requirements for a deep excavation can result in a fatal trench collapse. A missing Health and Safety File that does not record hidden service routes can result in a fatality when a future operative strikes an unidentified live cable. The likelihood of harm arising from CDM failures is directly proportional to the significance of the hazard that the failed document was intended to address and control. Principal contractors and clients who treat CDM documentation as a compliance exercise rather than a risk management tool are significantly more likely to have documentation failures that translate into foreseeable harm.}

\subsection{Step 4 --- Select and Implement Controls}
\paragraph{Controls for CDM compliance risk must operate from client appointment onward. The client must conduct a genuine, proportionate competence assessment of the organisations and individuals it appoints as Principal Designer and Principal Contractor, going beyond SSIP certificate verification to assess the appointee's actual knowledge and experience of the project type and complexity. The Principal Designer must be appointed in writing before the design process starts, not as an afterthought at planning application stage. The Construction Phase Plan must be prepared before construction starts, reviewed against the actual site conditions and pre-construction information, and updated throughout the construction phase as conditions change. The Health and Safety File must be maintained throughout the project by a designated owner, with a handover checklist that confirms completeness at practical completion. All persons on site must hold appropriate CSCS cards for their role; the site induction must verify card validity; and any operative whose card has expired or whose card level is not appropriate to their role must not be permitted to work until the card situation is resolved.}

\subsection{Step 5 --- Record and Review}
\paragraph{CDM documentation must be reviewed at each phase transition of the project and at any significant change in scope, contractor list, or site conditions. The Construction Phase Plan should carry a revision history showing each update and the reason for it. The Principal Designer's design risk register should be transferred to the Principal Contractor at the start of the construction phase, creating an unbroken chain of risk management information from design through construction. Monthly CDM compliance audits by the Principal Contractor's safety team should check: current CPP version against site activity; PCI completeness; F10 status; HSF progress; and CSCS card compliance across the current subcontractor list.}

\example{Five-Step Application}{A client appoints a Principal Designer for a \pounds4\,m commercial fit-out project. Step 1 identifies: the proposed PD is a small architectural practice with no documented CDM PD experience; PCI does not include asbestos refurbishment survey; F10 notification has not been submitted. Step 2 identifies: all construction phase workers at risk from inadequate design risk assessment; future tenants at risk from incomplete HSF. Step 3 rates risk from inadequate PCI (missing asbestos survey) as Very High given the 1980s construction date. Step 4 implements: client requests evidence of PD competence; R\&D asbestos survey commissioned before design stage commences; F10 submitted; PD appointment confirmed in writing with specific competence evidence reviewed. Step 5 specifies: CDM compliance audit at RIBA Stage 4, at start of construction phase, and at each significant contractor appointment.}

%---------------------------- SECTION ----------------------------
\section{Applying the hierarchy of controls}

\paragraph{The hierarchy of controls for CDM compliance risk is necessarily administrative and procedural in nature, given that CDM is itself a management framework rather than a physical hazard. However, the principles of the hierarchy are applicable: the highest-order control is the elimination of CDM failures by embedding compliance into the project management processes so that the risk of non-compliance is systematically reduced, not managed reactively.}

\begin{itemize}
  \bitem{Elimination}{Eliminate the risk of CDM failures by making CDM compliance a non-negotiable project gateway: no appointment without written competence evidence; no construction phase start without a complete CPP and F10 submission; no practical completion without HSF handover. Eliminate the risk of inadequate PCI by commissioning all required surveys (asbestos, structural, geotechnical, utility) as a planned, budgeted early project activity.}
  \bitem{Substitution}{Where a proposed appointment cannot demonstrate CDM competence, substitute a more capable appointee rather than proceeding with an inadequate duty holder. Substitute generic CPP templates with project-specific plans developed against the actual PCI.}
  \bitem{Engineering controls}{In the CDM context, ``engineering controls'' equate to systematic process controls: a CPP template that requires project-specific completion of each section and cannot be finalised without all mandatory elements; a HSF management platform that automatically flags missing information against the required content list; a project CSCS card management system that electronically verifies card validity against an up-to-date database.}
  \bitem{Administrative controls}{Implement CDM compliance audits at pre-start, monthly, and at phase transitions; maintain a CDM duty holder register for every project showing name, appointment date, and competence evidence reference; implement a CPP update procedure triggered by each significant change in site conditions or contractor composition; implement a PCI distribution log confirming receipt and acknowledgement by all designers and contractors.}
  \bitem{PPE}{In the CDM context, PPE represents the residual assurance measures: CDM compliance checklists used by site managers; HSF handover checklists used at practical completion; and post-completion CDM audits that verify the completeness of documentation before the HSF is archived.}
\end{itemize}

%---------------------------- SECTION ----------------------------
\section{Cadence of assessment and review triggers}

\paragraph{CDM documentation has a natural cadence aligned with the project lifecycle, but must also respond to reactive triggers arising from scope changes, personnel changes, and incidents. The cadence must be embedded in the project management system as non-negotiable review points.}

\begin{itemize}
  \bitem{Pre-appointment}{Before any CDM duty holder is appointed, competence evidence must be gathered and assessed. This is a mandatory pre-condition of appointment, not an afterthought.}
  \bitem{Pre-construction phase start}{The CPP must be complete, the F10 submitted, welfare facilities confirmed, and the PCI distributed to all contractors before a single operative arrives on site.}
  \bitem{Monthly during construction phase}{The CPP should be reviewed monthly against actual site activity and updated to reflect any changes in scope, contractor list, or significant risk. The HSF should be updated with as-built information as each phase of the work is completed, not left to a pre-handover scramble.}
  \bitem{At each significant scope change}{Any significant change to the scope of the project --- additional floors, revised structural approach, introduction of a new principal contractor --- must trigger a CDM review including F10 update if project notification thresholds have changed.}
  \bitem{At practical completion}{A CDM close-out review must confirm: HSF completeness against the content list; all O\&M manuals and as-built drawings received; F10 notification status current; and CPP archived for the prescribed retention period.}
  \bitem{After any CDM-related incident}{Any incident where a CDM documentation failure contributed to the incident (inadequate PCI, out-of-date CPP, missing competence record) must trigger an immediate review of all CDM documentation on the project and on any similar active projects.}
\end{itemize}

%---------------------------- SECTION ----------------------------
\section{Common failure modes}

\begin{itemize}
  \bitem{Principal Designer appointed in name only, with no substantive input}{The most common CDM 2015 failure across the UK construction industry is the appointment of a Principal Designer who fulfils the role on paper but does not carry out the substantive duties of coordinating health and safety in the pre-construction phase, managing designer contributions to the PCI, or producing a meaningful design risk register. This is typically an architectural practice that has agreed to be named as PD on the appointment form without understanding or resourcing the role.}
  \bitem{Construction Phase Plan is a template with no project-specific content}{CPPs that are copied from a generic template and submitted to the client with project name and dates substituted, but with no project-specific risk assessment, no site-specific welfare arrangements, and no reference to the specific hazards identified in the PCI, are a regulatory fiction. They satisfy the formal requirement to have a CPP but provide no meaningful risk management benefit.}
  \bitem{F10 notification not submitted or submitted late}{Projects that exceed the notification thresholds frequently proceed without F10 notification, either through ignorance of the threshold criteria or through informal delegation of the client's notification duty to the Principal Contractor, who then fails to complete it. The F10 is not merely administrative; it alerts the HSE to the project's existence and enables proactive inspection, which is a regulatory safeguard.}
  \bitem{Health and Safety File not produced or not handed over at completion}{The Health and Safety File is one of the most frequently missing CDM deliverables in the UK construction industry. Projects reach practical completion with incomplete O\&M manuals, missing as-built drawings, and absent hazardous material records, and the client accepts handover without insisting on HSF completeness. The building's future maintainers then work without essential safety information.}
  \bitem{CSCS card compliance not maintained during the construction phase}{CSCS card checks are conducted rigorously at the start of a project and during the first HSE or client audit, and then decline as the project progresses and the workforce changes. Operatives whose cards have expired, whose card level has changed, or who have joined the project mid-programme without a card check are a persistent compliance failure.}
\end{itemize}

\example{Failure Pattern}{A refurbishment contractor is appointed as both Principal Contractor and de facto Principal Designer (the client having failed to appoint a Principal Designer, leaving the duty to default to the contractor under CDM 2015 Regulation 5) on a \pounds1.2\,m commercial office refurbishment. The contractor produces a generic CPP that does not reference the R\&D asbestos survey that should have been commissioned as part of the PCI. Operatives begin drilling into partition walls during the fit-out phase. A routine HSE inspection discovers that the partitions contain chrysotile asbestos-insulating board. Work is immediately stopped; a Prohibition Notice is issued. An emergency R\&D survey reveals that 60\% of the partition walls contain ACMs. The licensed asbestos removal works add six weeks and \pounds180\,000 to the project. The contractor is prosecuted for failing to prepare an adequate CPP and failing to commission the pre-construction asbestos survey. The client is also prosecuted for failing to appoint a Principal Designer.}

%---------------------------- CHAPTER SUMMARY ----------------------------
\newpage
\section{Chapter Summary}

\begin{table}[H]
\centering
\begin{tabularx}{\textwidth}{>{\raggedright\arraybackslash}p{3.2cm}X}
\toprule
\textbf{Assessment Element} & \textbf{Operational Requirement} \\
\midrule
Legal/Professional Basis & CDM 2015; L153 ACOP; Health and Safety at Work etc.\ Act 1974. \\
Method & Five-step CDM compliance risk assessment; duty holder competence assessment; PCI, CPP, and HSF management process. \\
Controls & Written duty holder appointments with competence evidence; complete PCI before design starts; project-specific CPP before construction phase starts; F10 notification submitted on time; HSF completed and handed over at practical completion; CSCS card compliance maintained throughout. \\
Cadence & Pre-appointment; pre-construction phase start; monthly during construction phase; at each significant scope change; at practical completion; after any CDM-related incident. \\
Assurance & Monthly CDM compliance audit; CPP update log; HSF progress tracker; CSCS card compliance sweep; post-completion CDM documentation audit. \\
\bottomrule
\end{tabularx}
\end{table}

\begin{itemize}
  \bitem{Key takeaway 1}{CDM 2015 applies to all construction work in Great Britain; the domestic client provision does not exempt projects from CDM --- it transfers the client's duties to the contractor or Principal Contractor, who must discharge them fully.}
  \bitem{Key takeaway 2}{The Principal Designer must be a genuine CDM duty holder with the knowledge, skills, experience, and organisational capability to perform the pre-construction phase coordination role; appointment in name only without substantive input is both a CDM breach and a significant risk to the construction phase.}
  \bitem{Key takeaway 3}{The Construction Phase Plan must be project-specific, prepared before construction commences, and updated throughout the construction phase; a generic template does not constitute a CPP and provides no legal or operational protection.}
  \bitem{Key takeaway 4}{The Health and Safety File must be completed and handed to the client at practical completion; it is the principal contractor's legal obligation to ensure its completeness, and the client's legal obligation to keep it and make it available to future workers on the building.}
  \bitem{Key takeaway 5}{Competence assurance must go beyond CSCS card verification to assess whether the organisation or individual has the actual knowledge, skills, and experience to perform their CDM role on the specific project; this applies to all duty holder appointments from client to subcontractor.}
\end{itemize}

\barquote{In Chapter~\ref{chap:Subcontractor, Supply Chain and Procurement Risk Assessment}, we examine the risk assessment and management of the subcontractor and supply chain relationships that define the practical delivery of health and safety on the construction site.}
